% TOP_PROBLEM: {"id":30007715,"problem_number":"LOCAL-30007715","title":"Cost-effective climate adaptation","category_id":21,"category":{"id":21,"name":"economics","display_name":"Economics","description":"Open economic theory statements, published research agendas, and proposed scoped specifications; question provenance and historical priority are qualified per record.","slug":"economics","order_index":21,"created_at":"2026-10-08T00:00:00Z"},"status":"research_question","external_url":null,"legacy_ids":[],"published":false,"statement":"Choose a budget-feasible adaptation package for heat-exposed low-income districts that maximizes welfare under explicitly stated future climate distributions.","economics":{"original_id":204,"field":"Climate, environment and energy","kind":"design","formalization_basis":"adapted_research_question","source_status":"research_agenda","priority_status":"earliest_found","priority_note":"Earliest explicit statement located in this audit; no exhaustive priority search or first-ever claim. The mathematical specification below is adapted, not quoted from the source.","earliest_verified_reference":{"authors":["Tamma Carleton","Esther Duflo","Kelsey Jack","Guglielmo Zappalà"],"title":"The economics of climate adaptation: From academic insights to effective policy","year":2025,"url":"https://cepr.org/voxeu/columns/economics-climate-adaptation-academic-insights-effective-policy","doi":null,"locator":"Sections “The study of adaptation to inform adaptation policy” and “Future directions for research on adaptation”, 15 April 2025","evidence_summary":"The authors call for mechanism identification, future-climate extrapolation, and welfare analysis of adaptation barriers."},"current_reference":{"authors":["Tamma Carleton","Esther Duflo","Kelsey Jack","Guglielmo Zappalà"],"title":"The economics of climate adaptation: From academic insights to effective policy","year":2025,"url":"https://cepr.org/voxeu/columns/economics-climate-adaptation-academic-insights-effective-policy","doi":null,"locator":"Sections “The study of adaptation to inform adaptation policy” and “Future directions for research on adaptation”, 15 April 2025","evidence_summary":"The authors call for mechanism identification, future-climate extrapolation, and welfare analysis of adaptation barriers."},"formalization_note":"The source explicitly identifies the underlying gap or agenda. Population, horizon, interventions, estimands, and auxiliary assumptions are proposed operational choices, not a canonical source theorem. Partial later answers are compatible with this scoped question; the citation alone does not prove absence of every subsequent solution. Mathematical scope was checked independently; added definedness, feasibility, or identification conditions belong to this proposed formalization."},"statusQualification":"Published question or research agenda verified at the cited locator. The mathematical specification is adapted for this catalogue and includes additional assumptions; source publication does not establish first-ever priority or certify this exact model remains unsolved. The source explicitly identifies the underlying gap or agenda. Population, horizon, interventions, estimands, and auxiliary assumptions are proposed operational choices, not a canonical source theorem. Partial later answers are compatible with this scoped question; the citation alone does not prove absence of every subsequent solution. Mathematical scope was checked independently; added definedness, feasibility, or identification conditions belong to this proposed formalization.","statusReviewedAt":"2026-10-08"}
% FIRST_POSTED: 2026-10-08
% ENTRY_KIND: statement_only
% !TeX program = lualatex
\documentclass[11pt]{article}
\usepackage{fontspec}
\usepackage[a4paper,margin=25mm]{geometry}
\usepackage{amsmath,amssymb,mathtools}
\usepackage{xurl}
\usepackage[hidelinks]{hyperref}
\newcommand{\sref}[2]{\hyperlink{source:#1:#2}{[S#2]}}
\setlength{\emergencystretch}{3em}
\begin{document}
% problemId:30007715
\section*{Cost-effective climate adaptation}\label{30007715}
\subsection{Definitions and mathematical statement}
Consider low-income households in heat-exposed districts of one country and a thirty-year sequence of projected temperatures. Let \(a\in\mathcal A\) denote an implementable adaptation package, where the finite set \(\mathcal A\) includes cooling access, warning systems, and infrastructure upgrades. Define household discounted health and consumption utility \(U_i(a,c)\) under climate path \(c\), and discounted public resource cost \(K(a,c)\). The climate distribution \(P\) and welfare weights \(w_i>0\) must be stated before estimation. Include a zero-cost baseline package, assume finite expected utilities and costs, and let \(B\geq0\) be the budget. The target is
\[a^*(P,B)\in\arg\max_{a\in\mathcal A:E_P[K(a,c)]\leq B}E_P\!\left[\sum_i w_iU_i(a,c)\right].\]
Assume packages can be assigned at district level, but allow household responses and within-district spillovers. Estimate policy effects using credible intervention variation and identify which response parameters must be transported beyond observed weather. The open research agenda concerns welfare under future climates, rather than merely whether present-day cooling reduces heat harm. Compare rankings across plausible climate distributions, uptake costs, and later complementary investments. If available variation does not identify long-run responses, give explicit bounds or alternative scenario rankings. All package definitions, budgets, utility normalization, and transport restrictions are part of this proposed operationalization.

\par\textbf{Formulation and scope.} The source explicitly identifies the underlying gap or agenda. Population, horizon, interventions, estimands, and auxiliary assumptions are proposed operational choices, not a canonical source theorem. Partial later answers are compatible with this scoped question; the citation alone does not prove absence of every subsequent solution. Mathematical scope was checked independently; added definedness, feasibility, or identification conditions belong to this proposed formalization.

\par\textbf{Open-question provenance.} An explicit question or research agenda was verified at the source locator. This does not by itself certify that every later version remains unresolved \sref{30007715}{1}.

\par\textbf{Historical priority.} Earliest explicit statement located in this audit; no exhaustive priority search or first-ever claim. The mathematical specification below is adapted, not quoted from the source.

\subsection{Short English statement}
Choose a budget-feasible adaptation package for heat-exposed low-income districts that maximizes welfare under explicitly stated future climate distributions.

\subsection{Sources}
\begin{itemize}\raggedright
\item[\textnormal{[S1]}]\hypertarget{source:30007715:1}{} Tamma Carleton, Esther Duflo, Kelsey Jack, Guglielmo Zappalà. 2025. The economics of climate adaptation: From academic insights to effective policy. Sections “The study of adaptation to inform adaptation policy” and “Future directions for research on adaptation”, 15 April 2025.
\url{https://cepr.org/voxeu/columns/economics-climate-adaptation-academic-insights-effective-policy}.
{\small\textit{Earliest verified question/agenda reference; historical priority qualified below. The authors call for mechanism identification, future-climate extrapolation, and welfare analysis of adaptation barriers.}}
\end{itemize}
\end{document}
